\resumeSubHeadingListStart    
    \resumeSubheading
    {DataEngineer, AI Labs}{Jun 2022 -- Sep 2024}
    {}{중구, 서울}
    
    \resumeProjectHeading
    {
        \hypertarget{price_database_back}
            {
                \textbf{Price Database 서비스를 위한 데이터 파이프라인 개발}
            }
    }
    {Jan 2023 -- Oct 2023}
    \resumeItemListStart
        \resumeItem{개발 리소스를 아끼려고 AWS Aurora RDS 한 대로 데이터 파이프라인 구축}
        \resumeItem{AWS EMR에서 PySpark로 데이터 가공 · 적재}
        \resumeItem{분산 처리를 위한 파티셔닝 및 데이터 최적화}
        \resumeItem{데이터 플랫폼(Price Database/MLOps) 총괄 설계 및 구축으로 Series B 투자 유치에 핵심 기여}
        \resumeItem{\textbf{대용량 처리가 가능한 데이터 파이프라인을 주도해 구축}}
        \resumeItem{\hyperlink{price_database}{\underline{Project Detail >>}}}
    \resumeItemListEnd

    \resumeProjectHeading
    {
        \hypertarget{art_market_scraper_back}
            {
                \textbf{미술시장 데이터 스크래퍼 개발}
            }
    }
    {Jan 2023 -- Feb 2023}
    \resumeItemListStart
        \resumeItem{requests와 BeautifulSoup4로 데이터 스크래퍼 개발}
        \resumeItem{멀티스레딩으로 리소스를 빠르게 업로드하는 스크래퍼 개발}
        \resumeItem{LangChain 정보 추출 프롬프트로 정형 데이터 스크래핑 구현}
        \resumeItem{Swift로 작성하여 iOS 환경에서 동작하는 소셜 관계 수집 스크래퍼 개발}
        \resumeItem{boto3로 스크래퍼 인스턴스 50대를 동시에 관리}
        \resumeItem{WebSocket으로 관리 서버와 연결해 스크래퍼 인스턴스 제어}
        \resumeItem{\textbf{동적(Selenium, PlayWright), 정적(Public/Private API) 데이터 스크래퍼 개발 및 운용}}
        \resumeItem{\hyperlink{art_market_scraper}{\underline{Project Detail >>}}}
    \resumeItemListEnd

    \resumeProjectHeading
    {
        \hypertarget{scraper_server_back}
            {
                \textbf{데이터 스크래퍼 중앙 관리 서버 개발}
            }
    }
    {Jan 2023 -- Feb 2023}
    \resumeItemListStart
        \resumeItem{Flask로 스크래퍼를 중앙에서 관리하는 API 엔드포인트 개발}
        \resumeItem{WebSocket으로 스크래퍼와 연결해 수집 대상을 내려보내는 구조 구현}
        \resumeItem{SHA-256 해시로 로그인 자격 증명 전달}
        \resumeItem{Flask Blueprint로 엔드포인트 모듈화}
        \resumeItem{Nginx로 리버스 프록시 구성}
        \resumeItem{Route53과 Certbot으로 HTTPS 인증서 관리}
        \resumeItem{\textbf{스크래퍼를 한곳에서 관리하는 서버 개발 · 운영}}
        \resumeItem{\hyperlink{scraper_server}{\underline{Project Detail >>}}}
    \resumeItemListEnd

    \resumeProjectHeading
    {
        \hypertarget{BI_tool_back}
            {
                \textbf{BI 툴 개발}
            }
    }
    {Nov 2023 -- Nov 2023}
    \resumeItemListStart
        \resumeItem{Flask와 NextJS로 BI 툴 서버 개발}
        \resumeItem{소셜 관계 데이터를 전달해주는 API 개발}
        \resumeItem{Cytoscape와 SigmaJS로 소셜 관계를 방향 그래프로 보여주는 BI 페이지 개발}
        \resumeItem{데이터 라벨링을 위한 API 개발}
        \resumeItem{NextJS와 Carbon Design System으로 컴포넌트 개발}
        \resumeItem{Route53과 Certbot으로 HTTPS 인증서 관리}
        \resumeItem{Docker Compose로 EC2 인스턴스(Ubuntu 22.04)에 배포}
        \resumeItem{\textbf{BI 툴을 화면부터 인프라까지(Nginx · NextJS · Flask) 직접 개발 · 운영}}
        \resumeItem{\hyperlink{BI_tool}{\underline{Project Detail >>}}}
    \resumeItemListEnd

    \resumeProjectHeading
    {
        \hypertarget{airflow_back}
            {
                \textbf{Airflow 서비스 배포 \& DAG 작성}
            }
    }
    {Jan 2023 -- Feb 2024}
    \resumeItemListStart
        \resumeItem{Airflow PythonOperator로 관리 서버와 통신하는 스크래핑 스케줄링 DAG 작성}
        \resumeItem{boto3로 AWS EMR을 띄워 데이터 정제 · 적재를 돌리는 DAG 작성}
        \resumeItem{Sequential Executor로 DAG 운영}
        \resumeItem{Route53과 Certbot으로 HTTPS 인증서 관리}
        \resumeItem{Docker Compose로 EC2 인스턴스(Ubuntu 22.04)에 서비스 배포}
        \resumeItem{\textbf{Airflow DAG로 스크래핑을 주기 실행하고 PySpark로 데이터 정제}}
        \resumeItem{\hyperlink{airflow}{\underline{Project Detail >>}}}
    \resumeItemListEnd

    \resumeProjectHeading
    {
        \hypertarget{removing_duplicates_back}
            {
                \textbf{데이터 중복 제거 기능 개발}
            }
    }
    {Nov 2023 -- Feb 2024}
    \resumeItemListStart
        \resumeItem{ML 리서치 엔지니어와 함께 Faiss-CPU로 IVFFlat 인덱스를 만들고, 임베딩 벡터 거리를 계산해 중복 제거 시스템 개발}
        \resumeItem{Facebook SSCD 모델로 이미지 임베딩 벡터 추출}
        \resumeItem{OpenAI 텍스트 임베딩 모델 Ada로 텍스트 임베딩 벡터 추출}
        \resumeItem{PostgreSQL pgvector 0.5.0 확장으로 인덱스 생성 오버헤드를 50\% 줄임}
        \resumeItem{pgvector의 IVFFlat · HNSW 인덱스 생성}
        \resumeItem{기존 numpy.array를 pgvector::vector로 전환}
        \resumeItem{\textbf{신뢰도 높은 데이터를 위한 Embedding 기반 데이터 정제}}
        \resumeItem{\hyperlink{removing_duplicates}{\underline{Project Detail >>}}}
    \resumeItemListEnd

    \resumeProjectHeading
    {
        \hypertarget{data_warehouse_back}
            {
                \textbf{효율적인 데이터 관리를 위해 데이터 웨어하우스 스키마 설계}
            }
    }
    {Feb 2024 -- Sep 2024}
    \resumeItemListStart
        \resumeItem{여러 소스에서 들어오며 생기는 중복 데이터를 관리하려고 SSoT 테이블을 포함한 데이터 웨어하우스를 DBDiagram으로 설계}
        \resumeItem{RDS 하나가 데이터 소스와 웨어하우스를 겸하며 커진 비용을 줄이려고, S3를 데이터 레이크로 쓰는 스키마로 재설계}
        \resumeItem{DataSource · DataWarehouse · DataMart · SSoT 스키마와 테이블 약 50개 설계}
        \resumeItem{\textbf{비용 절감 및 신뢰도 높은 데이터 관리를 위한 데이터 파이프라인 구조 재설계 주도}}
        \resumeItem{\hyperlink{data_warehouse}{\underline{Project Detail >>}}}
    \resumeItemListEnd

    \resumeProjectHeading
    {\textbf{AWS Service 관리}}{}
    \resumeItemListStart
        \resumeItem{IAM Role 기반 유저 별 서비스 접근 제어}
        \resumeItem{비용을 줄이려고 EC2 인스턴스 타입을 x86에서 arm64로 전환}
        \resumeItem{읽기 · 쓰기 인스턴스를 분리해 Aurora RDS PostgreSQL 운영}
        \resumeItem{매니지드 서비스와 EC2 직접 운영의 비용을 따져 골라 씀}
        \resumeItem{\textbf{리소스별로 필요한 만큼만 권한 부여}}
    \resumeItemListEnd
\resumeSubHeadingListEnd